\newpage
\section{Stand der Forschung}

\subsection{LLM-gestützte Testgenerierung}

Die automatisierte Testgenerierung mit \acp{LLM} ist eines der Anwendungsfelder des in Kapitel~2.1.2 beschriebenen Einsatzes von \acp{LLM} im Software Engineering. \textcite{celik_review_2025} fassen die bestehenden Ansätze in einer Literaturübersicht zusammen und unterscheiden zwischen vier methodischen Kategorien: Prompt Engineering, feedback-getriebene Verfahren, Fine-Tuning sowie hybride Kombinationen aus \acp{LLM} und traditionellen Testwerkzeugen. Diese Arbeit gehört zur ersten Kategorie, da sie das Modellverhalten über die Eingabe steuert. Die ausgewerteten Studien betreffen überwiegend Java und Python, weil Testwerkzeuge und Benchmarks vor allem für diese Sprachen verfügbar sind (\cite{celik_review_2025}).

Dass eine natürlichsprachliche Beschreibung als fachliche Eingabe ausreichen kann, zeigen \textcite{plein_automatic_2023} in einer Machbarkeitsstudie. Sie lassen ChatGPT aus natürlichsprachlichen Bug Reports ausführbare Testfälle erzeugen und bewerten das Ergebnis anhand der Ausführbarkeit, Validität und Relevanz. Bei fünf Generierungsversuchen je Bug Report entsteht für 50\,\% der Eingaben ein ausführbarer Test, für etwa 30\,\% ein Test, der den berichteten Fehler tatsächlich reproduziert, und für etwa 9\,\% ein Test, der zusätzlich die korrigierte Programmversion passiert (\cite{plein_automatic_2023}). Die häufigsten Ursachen der nicht ausführbaren Ausgaben sind fehlende Imports, doppelte Methodennamen und veraltete Funktionsaufrufe. Für die Fragestellung dieser Arbeit ist die Studie allerdings nur bedingt aussagekräftig, da sie Java-Projekte betrachtet und die generierten JUnit-Testfälle an die bestehende Testsuite des jeweiligen Projektes anhängt. Eine Benutzeroberfläche kommt nicht vor, weshalb sich die Ergebnisse nur bedingt auf \ac{E2E}-Tests übertragen lassen.

Für den \ac{E2E}-Testbereich liegt eine industrielle Fallstudie von \textcite{ferreira_acceptance_2025} vor. Bei dem dort vorgestellten Verfahren erzeugt das Werkzeug AutoUAT mit GPT-4 Turbo zunächst aus User Stories Testszenarien im Gherkin-Format. Das zweite Werkzeug, Test-Flow, generiert daraus anschließend ausführbare Tests für das Testframework Cypress. Als Kontext erhält das Modell neben der User Story auch den \ac{HTML}-Code der betroffenen Seiten, aus dem Style- und Script-Elemente vorab entfernt werden (\cite{ferreira_acceptance_2025}).

Von 65 abgegebenen Nutzerrückmeldungen zu AutoUAT bewerteten 95\,\% das generierte Szenario als hilfreich (\cite{ferreira_acceptance_2025}). Test-Flow wurde an 13 User Stories mit insgesamt 50 Gherkin-Szenarien gemessen. Von den 50 generierten Testfällen waren 60\,\% unverändert einsetzbar und 8\,\% nach kleinen Korrekturen. 24\,\% ließen sich erst verwenden, nachdem die Eingabe um den fehlenden Kontext ergänzt wurde, und 8\,\% mussten wegen schwerer Fehler verworfen werden. Fehlender Kontext war mit 12 von 20 fehlerhaften Testfällen die häufigste Fehlerursache (\cite{ferreira_acceptance_2025}). Die verwendeten User Stories verteilen sich auf mehrere Seitentypen sowie formular- und navigationsbasierte Interaktionen eines E-Commerce-Produkts.

Einen anderen Weg zur Kontextgewinnung beschreiben \textcite{alian_feature-driven_2024} mit AutoE2E. Das Verfahren leitet die Features einer Webanwendung aus deren beobachtetem Zustand ab, statt sie einer Anforderungsbeschreibung zu entnehmen. Eine Breitensuche crawlt den Zustandsraum einer Anwendung, bis keine neuen Zustände mehr gefunden werden, und ein \ac{LLM} erzeugt daraus ausführbare \ac{E2E}-Testfälle. Auf dem von den Autoren erstellten Benchmark E2EBench erreicht AutoE2E eine durchschnittliche Feature-Abdeckung von 79\,\%, gegenüber 12\,\% beim klassischen Verfahren mit Crawljax (\cite{alian_feature-driven_2024}). Nach eigener Angabe ist AutoE2E das erste Verfahren, das feature-getriebene \ac{E2E}-Tests autonom generiert.

Gemeinsam zeigen die betrachteten Arbeiten, dass sich ausführbare Tests für Webanwendungen aus natürlichsprachlichen Eingaben erzeugen lassen (\cite{plein_automatic_2023}; \cite{ferreira_acceptance_2025}; \cite{alian_feature-driven_2024}). Alle Ansätze nutzen dabei jeweils eine andere Form von Anwendungs- oder \ac{UI}-Kontext. Allerdings untersucht keine der Arbeiten, wie sich Form und Umfang des bereitgestellten \ac{UI}-Kontexts auf die Qualität der generierten Tests auswirken.

\subsection{Forschungslücke und Einordnung dieser Arbeit}

Ein Bezug zu WebGIS-Anwendungen ist in keiner der genannten Arbeiten vorhanden. Nach eigener Recherche liegen keine wissenschaftlichen Arbeiten zur \ac{LLM}-gestützten Testgenerierung für WebGIS-Anwendungen mit Karteninteraktionen vor. Auch für \ac{OPT} als Framework existieren keine Untersuchungen. Die in Kapitel~3.1 betrachteten Ansätze setzen voraus, dass die zu prüfenden Elemente über das \ac{DOM} erreichbar und ansprechbar sind. Für die Karte gilt diese Annahme nicht (vgl. Kapitel~4.2), sodass offen ist, ob sich ihre Ergebnisse auf diesen Anwendungsfall übertragen lassen. Der Einfluss des \ac{UI}-Kontexts wiederum ist durch das In-Context Learning erklärbar (vgl. Kapitel~2.1.1), wurde in der Literatur bisher aber nur als Nebenbefund berichtet und nicht kontrolliert untersucht.

Diese Arbeit setzt an beiden Lücken an. Sie knüpft an den Workflow von der Anforderung zum Test an (\cite{plein_automatic_2023}; \cite{ferreira_acceptance_2025}) und grenzt sich von editorbasierten Werkzeugen ab, die aus dem umgebenden Quellcode generieren. Ihr Beitrag besteht darin, den \ac{UI}-Kontext als kontrollierte Variable zu untersuchen und erstmals auf WebGIS-Anwendungen einschließlich Karteninteraktionen zu beziehen.
